\section{Prime returns from a positive graph}
\label{sec:prime-returns}
Repeated prime shifts can be generated by a unitary return channel.
The construction fixes their full geometric series and displays the
contact required by positive output power.

\subsection{A nonnegative graph Hamiltonian and its scattering matrix}
For a prime $p$, set $q=p^{-1/2}$ and $d=\log p$.
Use two half-line leads and a finite stub of length $d/2$ with a
Neumann endpoint. At the common vertex choose the real symmetric
orthogonal matrix
\[
O_q=
\begin{pmatrix}
1-q&\sqrt q&\sqrt{q(1-q)}\\
\sqrt q&0&-\sqrt{1-q}\\
\sqrt{q(1-q)}&-\sqrt{1-q}&q
\end{pmatrix},\qquad P_\pm=(I\pm O_q)/2.
\]
The quadratic form $\sum_{\text{edges}}\int|\psi'|^2$ on edgewise
$H^1$ functions with $P_-\psi(0)=0$ is closed and nonnegative.
Its self-adjoint Laplacian has vertex conditions
$P_-\psi(0)=0$, $P_+\psi'(0)=0$, with outgoing derivatives.
This is the standard self-adjoint graph-scattering framework
\cite{KostrykinSchrader}, specialized to an explicit coupler.

Write $O_q$ in blocks
$\left(\begin{smallmatrix}D_q&B_q^T\\B_q&q\end{smallmatrix}\right)$,
where
$D_q=\left(\begin{smallmatrix}1-q&\sqrt q\\\sqrt q&0\end{smallmatrix}\right)$
and $B_q=\sqrt{1-q}(\sqrt q,-1)$.
The stub round trip at wave number $k$ multiplies by $z=e^{ikd}$.
For lead input $a$, stub amplitude $v$ and output $b$,
$v=B_qa+qzv$, $b=D_qa+B_q^Tzv$. Elimination gives
\begin{equation}\label{prime:scattering}
V_q(z)=\frac1{1-qz}
\begin{pmatrix}
1-q&\sqrt q(1-z)\\
\sqrt q(1-z)&(1-q)z
\end{pmatrix},\qquad
\det V_q(z)=\frac{z-q}{1-qz}.
\end{equation}
The matrix is unitary for $|z|=1$. Its complementary outputs for
one incident lead are
$K_0=(1-q)/(1-qz)$ and
$K_1=\sqrt q(1-z)/(1-qz)$, with $|K_0|^2+|K_1|^2=1$.

\subsection{The prime source and its compulsory contact}
Feed the Fourier input into the complementary port with prescribed
scale $\sqrt{d(1+q)/(1-q)}$, and attach the unit cap $\chi$ if
desired. Equivalently, with $U=U_d$, the bounded whole-line source is
\begin{equation}\label{prime:source}
\mathcal D_pF=
\sqrt{\frac{dq(1+q)}{1-q}}\,
 (I-U)(I-qU)^{-1}F.
\end{equation}
The inverse is its norm-convergent Neumann series. Two graph copies
fed by $\widehat F(k)$ and $\widehat F(-k)$ accommodate the two
Fourier signs with nonnegative graph energy $k^2$. The arithmetic
Fourier variable is a wave number in this construction, not a
physical evolution time. The Hilbert adjoint of
\eqref{prime:source} follows by replacing $U$ by $U^*$ and reversing
the factors, and the interval adjoint includes $E_L^*$.

The Poisson kernel
\[
\Pi_q(z)=\frac{1-q^2}{|1-qz|^2}
        =1+\sum_{m\ge1}q^m(z^m+z^{-m})
\]
gives the exact output identity
\begin{equation}\label{prime:gram}
\mathcal D_p^*\mathcal D_p
=\frac{d(1+q)}{1-q}I-d\Pi_q(U)
=c_pI-d\sum_{m\ge1}q^m(U^m+U^{-m}),
\quad c_p=\frac{2dq}{1-q}.
\end{equation}
Moreover $\norm{\mathcal D_p}^2=4dq/(1-q^2)$.
Compression deletes precisely those terms with $md\ge L$.
Thus the channel gives every repetition with the arithmetic weight,
but also a positive contact $c_p$. The incident scale was prescribed
to obtain those weights; vertex unitarity alone does not choose
that scale.

\begin{proposition}\label{prime:contact}
For a finite prime set $\mathcal S$, a stationary whole-line Gram
operator with off-diagonal part
$-\sum_{p\in\mathcal S}d_p\sum_{m\ge1}q_p^m(U_{md_p}+U_{-md_p})$
and constant diagonal $cI$ is nonnegative if and only if
$c\ge c_{\mathcal S}:=\sum_{p\in\mathcal S}c_p$.
\end{proposition}
\begin{proof}
Each cosine in the multiplier is at most one, which proves
sufficiency. At frequency zero they are simultaneously one.
Fourier packets concentrated at zero prove necessity. The argument
applies to arbitrary coherent realizations of the same stationary
Gram operator.
\end{proof}
This lower bound is for the whole-line stationary model. At fixed
support a smaller contact can suffice. For one prime, decompose
$I_L$ into translation chains modulo $d$. The largest chain has
$n=\lceil L/d\rceil$ points, almost everywhere when this length
occurs, and the correlation matrix on such a chain is
$R_n(q)=(q^{|i-j|})_{i,j=1}^n$. The exact interval contact is
\begin{equation}\label{prime:finite-contact}
c_{p,L}^{\min}=d\bigl(\lambda_{\max}(R_n(q))-1\bigr).
\end{equation}
It is zero for $L\le d$ and $dq$ for $d<L\le2d$.
Testing the long chain on the constant vector, together with the
row-sum upper bound, shows that it tends to $c_p$ as $L\to\infty$.

\subsection{Phase derivative and the Euler logarithm}
For $Z(s)=(1-p^{-s})^{-1}$ on $\re s=1/2$,
\[
2\re\,\partial_s\log Z(s)=d-d\Pi_q(e^{-i\,d\,\im s}).
\]
Its contact is therefore different from the output-power contact
in \eqref{prime:gram}. In contrast, the scattering phase satisfies
\[
\partial_k\arg\det V_q(e^{ikd})=d\Pi_q(e^{ikd})>0,\qquad
-iV_q^\dagger\partial_kV_q
 =\frac{d}{|1-qe^{ikd}|^2}B_q^TB_q\ge0.
\]
The trace is the phase derivative, and differentiation with
respect to energy $k^2$ introduces $1/(2k)$ for $k>0$.
These formulas distinguish Euler logarithmic data, positive
output power and positive delay; none identifies them with the
complete arithmetic norm.

For a fixed interval it is sufficient to take
$\mathcal S_L=\{p:\log p<L\}$. An unrenormalized direct sum over
all primes fails even on a nonzero compactly supported input.
For the inactive primes all translates are disjoint from the
input, so their output norm is exactly $c_p\norm f^2$; the sum
of $c_p$ diverges (already it dominates a constant multiple of
$\sum_p1/p$). The finite-prime graph therefore does not itself
define an all-prime positive limit.
